% --- 摘要章节 (详细版) ---
% --- 文件路径: frontmatter/01-abstract.tex ---


% 设定页眉标记
\markboth{摘要}{摘要}


\begin{abstract}
    
本文\cite{developers2022tensorflow}聚焦于乡村振兴背景下现代农业的精准规划问题。针对某村庄在有限的异构性耕地资源（含露天耕地与大棚）和复杂的农艺要求（如轮作、防重茬）下，如何制定 2024--2030 年的最优种植方案，我们构建了一套由浅入深、层层递进的优化模型体系，旨在最大化长期经济总收益并系统性应对从稳定到波动的市场风险。

针对\textbf{问题一}，我们在确定性情景下建立以七年期总净收益最大化为目标的\textbf{多周期混合整数规划（Multi-period MIP）模型}，核心连续变量为年度-季节-地块-作物的种植面积，并通过二进制变量控制地块分散度。约束体系包括：(1) 土地面积约束；(2) 作物-土地适宜性；(3) 禁止重茬；(4) 豆类三年一次轮作等。以 2023 年数据为初始条件，使用 Gurobi 求解得到“超产滞销”与“超产降价 50\% 销售”两情景的最优可执行种植计划，详表填入 \texttt{result1\_1.xlsx} 与 \texttt{result1\_2.xlsx}。

针对\textbf{问题二}，考虑销量、亩产、成本及部分售价区间波动，引入\textbf{鲁棒优化（Robust Optimization）}并构造盒式不确定集，将目标转为最大化“最坏情景”保底净收益。通过推导原模型的鲁棒对应形式，得到单一静态七年方案；虽在乐观情景下非全局最优，但确保任意参数组合下的可行性与收益下界，结果汇总于 \texttt{result2.xlsx}。

针对\textbf{问题三}，我们计划（此处待补充完整内容）进一步引入： (1) 作物间替代与互补性；(2) 价格-供给弹性反馈；(3) 多阶段信息更新与再优化（滚动优化 / 自适应策略）；(4) 风险偏好（CVaR 或下偏风险）刻画；(5) 资源扩张或基础设施投资情景分析等。\textbf{此段为占位符，请补写实际推导、模型形式与结论。}

\vspace{0.5\baselineskip}
\noindent 关键词：多周期混合整数规划；鲁棒优化；轮作约束；农业可持续；风险管理

\end{abstract}

